% !PRINTED 610
\weberpage{610}
%<<CONT

1. \gesperrt{If $\Omega$ is a quadratic field with the fundamental number $\varDelta$, and $\varphi$ an integral functional in this field, then a quadratic irrational number belonging to the discriminant $\varDelta$}

\weq{17}{\omega = \frac{b + \sqrt{\varDelta}}{2c}}

\gesperrt{and a positive integral rational number $e$ can be so determined that}

\weq{18}{ec(t_1 + \omega t_2)}

\gesperrt{becomes a basis form of $\varphi$. Conversely, if $\omega$ is any quadratic irrational number belonging to the discriminant $\varDelta$, the form (18) is always a basis form of a functional $\varphi$.}

\begin{center}
§. 167.

\gesperrt{Equivalent Functionals and Equivalent Numbers.}
\end{center}

Now if $\omega, \omega_1$ are any two quadratic irrational numbers belonging to the discriminant $\varDelta$, we obtain two basis forms

\weq{1}{\begin{aligned} \varphi &= ec(t_1+\omega t_2)\\ \varphi_1 &= e_1c_1(t_1+\omega_1 t_2). \end{aligned}}

If now the two numbers $\omega, \omega_1$ are equivalent in the sense in which we stipulated this concept in \S.\,120 of the first Band, then

\weq{2}{\omega_1 = \frac{p\omega+q}{r\omega+s},}

in which $p,q,r,s$ are integral rational numbers which satisfy the condition

\weq{3}{ps - qr = \pm 1.}

If we then make the substitution (2) in $\varphi_1$, it follows that

\[ (r\omega+s)\varphi_1 = e_1c_1[st_1+qt_2+\omega(rt_1+pt_2)], \]

or, if we set

\weq{4}{u_1 = st_1+qt_2,\quad u_2 = rt_1+pt_2,}

then

\[ (r\omega+s)\varphi_1 = e_1c_1(u_1+\omega u_2). \]

Now $ec(u_1+\omega u_2)$ is associated with $\varphi$, because (4) is an
% German breaks mid-word at the page end: "eine ganz-" (word completed on the next page)
%CONT>>
